% !TEX root = ../main.tex

\section{Multi-dimensional Models and the Fundamental Theorems}

\begin{enumerate}
    \item \textbf{Multi-dimensional Wiener Process:} Many independent Wiener processes can drive the randomness in stock prices. A $d$-dimensional (vector valued) Wiener process is
    \[\vec{\mathbb{W}}(t) = \bigl(\mathbb{W}_1(t), \mathbb{W}_2(t), \dots, \mathbb{W}_d(t)\bigr),\]
    where
    \begin{enumerate}
        \item each $\mathbb{W}_i$ is a Wiener process,
        \item $\mathbb{W}_i$ and $\mathbb{W}_j$, $i \neq j$, are independent random variables, i.e.\ $\mathbb{P}\{\mathbb{W}_i(t) \in B, \mathbb{W}_j(t) \in B'\} = \mathbb{P}\{\mathbb{W}_i(t) \in B\}\mathbb{P}\{\mathbb{W}_j(t) \in B'\}$, equivalently $\sigma(\mathbb{W}_i)$ and $\sigma(\mathbb{W}_j)$ are independent.
    \end{enumerate}

    \item \textbf{Filtration of $\vec{\mathbb{W}}$:} $\{\mathcal{F}(t)\}_{0 \leq t \leq T}$ with $\mathcal{F}(s) \subseteq \mathcal{F}(t)$ for $s \leq t$, generated by the full collection of the $\mathbb{W}_i$, such that
    \begin{enumerate}
        \item $\vec{\mathbb{W}}(t)$ is $\mathcal{F}(t)$-measurable,
        \item $\vec{\mathbb{W}}(t) - \vec{\mathbb{W}}(s)$ is independent of the information in $\mathcal{F}(s)$.
    \end{enumerate}

    \item \textbf{Quadratic variations:} For each $i$, $d\mathbb{W}_i\, d\mathbb{W}_i = dt$, i.e.\ $[\mathbb{W}_i, \mathbb{W}_i](t) = t$. For $i \neq j$, $d\mathbb{W}_i\, d\mathbb{W}_j = 0$, i.e.\ $[\mathbb{W}_i, \mathbb{W}_j](t) = 0$ (due to independence).

    \item \textbf{Ito process (single, driven by $d$ Wiener processes):}
    \[d\mathbb{X}(t) = \Theta(t)\, dt + \vec{\sigma}(t)\cdot d\vec{\mathbb{W}}(t),\]
    with $\Theta$ and $\vec{\sigma}$ adapted.

    \item \textbf{Collection of Ito processes:} For $\alpha = 1, 2, \dots, n$,
    \[d\mathbb{X}_\alpha(t) = \Theta_\alpha(t)\, dt + \sum_{i=1}^{d}\sigma_{\alpha i}(t)\, d\mathbb{W}_i(t).\]
    For instance, with $n = 3$ processes and $d = 2$ Wiener processes,
    \[\begin{pmatrix}d\mathbb{X}_1 \\ d\mathbb{X}_2 \\ d\mathbb{X}_3\end{pmatrix} = \begin{pmatrix}\Theta_1 \\ \Theta_2 \\ \Theta_3\end{pmatrix}dt + \begin{pmatrix}\sigma_{11} & \sigma_{12} \\ \sigma_{21} & \sigma_{22} \\ \sigma_{31} & \sigma_{32}\end{pmatrix}\begin{pmatrix}d\mathbb{W}_1 \\ d\mathbb{W}_2\end{pmatrix}.\]
    In general $d\mathbb{X}$ and $\Theta$ are $n \times 1$ vectors, $d\vec{\mathbb{W}}$ is $d \times 1$, and $\sigma$ is an $n \times d$ matrix. $\sigma_{\alpha i}$ is the ``volatility'' of $\mathbb{X}_\alpha$ corresponding to the Wiener process $\mathbb{W}_i$.

    \item \textbf{Quadratic variation of a vector valued Ito process.} Retaining only the $d\mathbb{W}\, d\mathbb{W}$ terms,
    \[d\mathbb{X}_\alpha\, d\mathbb{X}_\alpha = \sum_{i=1}^{d}\sum_{j=1}^{d}\sigma_{\alpha i}\sigma_{\alpha j}\, d\mathbb{W}_i\, d\mathbb{W}_j = \sum_{i=1}^{d}\sigma_{\alpha i}^2\, dt\]
    since the terms with $i \neq j$ vanish. Similarly, the cross variation is
    \[\boxed{d\mathbb{X}_\alpha\, d\mathbb{X}_\alpha = \left(\sum_{i=1}^{d}\sigma_{\alpha i}^2\right)dt, \qquad d\mathbb{X}_\alpha\, d\mathbb{X}_\beta = \left(\sum_{i=1}^{d}\sigma_{\alpha i}\sigma_{\beta i}\right)dt}\]

    \item \textbf{Ito-Doeblin formula (multi-dimensional).} In one variable,
    \[df(t, x) = \partial_t f\, dt + \partial_x f\, dx + \frac{1}{2}\partial_x^2 f\, dx\, dx.\]
    For two variables,
    \[df(t, x, y) = \partial_t f\, dt + \partial_x f\, dx + \partial_y f\, dy + \frac{1}{2}\partial_x^2 f\, dx\, dx + \frac{1}{2}\partial_y^2 f\, dy\, dy + \partial_x\partial_y f\, dx\, dy.\]
    Note that there is no $\frac{1}{2}$ on the mixed term: the two cross terms $\frac{1}{2}\partial_x\partial_y f\, dx\, dy + \frac{1}{2}\partial_y\partial_x f\, dy\, dx$ combine because $\partial_x\partial_y f = \partial_y\partial_x f$ (true for ``nice'' functions). Likewise for three variables $x, y, z$, there are three squared terms with $\frac{1}{2}$ and three mixed terms $\partial_x\partial_y f\, dx\, dy + \partial_y\partial_z f\, dy\, dz + \partial_z\partial_x f\, dz\, dx$ without it.

    \item \textbf{Girsanov's Theorem ($d$-dimensional).} Let $\vec{\mathbb{W}}(t)$ be a $d$-dimensional Wiener process with filtration $\{\mathcal{F}(t)\}_{0 \leq t \leq T}$ and $\vec{\Theta}(t)$ a $d$-dimensional adapted process. Define
    \[\mathbb{Z}(t) = \exp\!\left(-\int_0^t\vec{\Theta}(u)\cdot d\vec{\mathbb{W}}(u) - \frac{1}{2}\int_0^t\vec{\Theta}(u)\cdot\vec{\Theta}(u)\, du\right).\]
    Then
    \begin{enumerate}
        \item $\mathbb{Z}(T)$ is an RN derivative ($\mathbb{P}\{\mathbb{Z}(T) > 0\} = 1$, $\mathbb{E}[\mathbb{Z}(T)] = 1$),
        \item $\mathbb{Z}(t)$ is an RN derivative process, with $\tilde{\mathbb{P}}(A) = \int_A\mathbb{Z}(T)\, d\mathbb{P}$,
        \item $\tilde{\vec{\mathbb{W}}}(t) = \vec{\mathbb{W}}(t) + \int_0^t\vec{\Theta}(u)\, du$ is a $d$-dimensional Wiener process under $\tilde{\mathbb{P}}$.
    \end{enumerate}

    \item \textbf{Multi-dimensional Market Model.} A money market with risk-free rate $\mathbb{R}(t)$ and $m$ stocks $\mathbb{S}_I$, $I = 1, \dots, m$, driven by $d$ Wiener processes:
    \[d\mathbb{S}_I(t) = \alpha_I(t)\mathbb{S}_I(t)\, dt + \mathbb{S}_I(t)\sum_{j=1}^{d}\sigma_{Ij}(t)\, d\mathbb{W}_j(t),\]
    where $\sigma_{Ij}$ are the volatility coefficients. The discount process is $d\mathbb{D}(t) = -\mathbb{D}(t)\mathbb{R}(t)\, dt$. We want to find a fair price for a derivative security written on these $m$ stocks for each $t \in [0, T]$, knowing the payoff at time $T$.

    \item \textbf{Hedging portfolio and wealth process:} With $\Delta_I(t)$ shares of $\mathbb{S}_I$ and the rest in the money market,
    \[d\mathbb{X}(t) = \sum_{I=1}^{m}\Delta_I(t)\, d\mathbb{S}_I(t) + \mathbb{R}(t)\left(\mathbb{X}(t) - \sum_{I=1}^{m}\Delta_I(t)\mathbb{S}_I(t)\right)dt.\]

    \item \textbf{Is there a risk-neutral measure?} We need a measure $\tilde{\mathbb{P}}$ such that every discounted stock price $\mathbb{D}(t)\mathbb{S}_I(t)$ is a martingale, i.e.\ $\tilde{\mathbb{E}}[\mathbb{D}(t)\mathbb{S}_I(t) \mid \mathcal{F}(s)] = \mathbb{D}(s)\mathbb{S}_I(s)$ for $s < t$ (equivalently $\mathbb{E}[\mathbb{Z}\mathbb{D}(t)\mathbb{S}_I(t) \mid \mathcal{F}(s)] = \mathbb{Z}(s)\mathbb{D}(s)\mathbb{S}_I(s)$). Requirement: the discounted stock has no $dt$ term under $\tilde{\mathbb{P}}$,
    \[d(\mathbb{D}\mathbb{S}_I) = \mathbb{D}\mathbb{S}_I\sum_{j=1}^{d}\sigma_{Ij}\, d\tilde{\mathbb{W}}_j.\]
    But what we know (by Ito's product formula) is
    \[d(\mathbb{D}\mathbb{S}_I) = (\alpha_I - \mathbb{R})\mathbb{D}\mathbb{S}_I\, dt + \mathbb{D}\mathbb{S}_I\sum_{j=1}^{d}\sigma_{Ij}\, d\mathbb{W}_j.\]
    Assume $d\tilde{\mathbb{W}}_j = d\mathbb{W}_j + \Theta_j\, dt$. Substituting, and comparing the $dt$ terms with the requirement, we get the \textbf{market price of risk equations}:
    \[\boxed{\alpha_I(t) - \mathbb{R}(t) = \sum_{j=1}^{d}\sigma_{Ij}(t)\Theta_j(t), \qquad I = 1, \dots, m} \tag{$\ast$}\]
    or in matrix form $(\alpha - \mathbb{R}\mathbf{1}) = \sigma\vec{\Theta}$, with $\sigma$ the $m \times d$ matrix.

    \item \textbf{Solutions of $(\ast)$:} These are $m$ linear equations in the $d$ unknowns $\Theta_j$. There are three possibilities:
    \begin{enumerate}
        \item \textit{No solution:} no $\tilde{\mathbb{P}}$ under which the $\mathbb{D}\mathbb{S}_I$ are martingales; this leads to arbitrage.
        \item \textit{Unique solution.}
        \item \textit{Infinitely many solutions.}
    \end{enumerate}
    If there is at least one solution, a risk-neutral measure $\tilde{\mathbb{P}}$ exists (via Girsanov with that $\vec{\Theta}$).

    \item \textbf{Arbitrage:} A wealth process $\mathbb{X}(t)$ such that
    \begin{enumerate}
        \item $\mathbb{X}(0) = 0$ (start with zero wealth),
        \item for some $T > 0$, $\mathbb{P}\{\mathbb{X}(T) \geq 0\} = 1$ (almost surely there is no loss of wealth),
        \item $\mathbb{P}\{\mathbb{X}(T) > 0\} > 0$ (some positive probability of gaining wealth).
    \end{enumerate}

    \item \textbf{First Fundamental Theorem of Asset Pricing:} If the market price of risk equations $(\ast)$ have at least one solution (equivalently, if there is a risk-neutral measure $\tilde{\mathbb{P}}$), then there is no arbitrage.
    \begin{proof}
        Suppose an arbitrage portfolio $\mathbb{X}(t)$ made of shares in the $\mathbb{S}_I$ and a money market account exists. Since $\tilde{\mathbb{P}}$ exists, $\mathbb{D}(t)\mathbb{X}(t)$ is a martingale under it, so
        \[\tilde{\mathbb{E}}[\mathbb{D}(T)\mathbb{X}(T)] = \mathbb{D}(0)\mathbb{X}(0) = 0. \tag{A}\]
        From $\mathbb{P}\{\mathbb{X}(T) \geq 0\} = 1$ and $\mathbb{D}(T) > 0$ we get $\mathbb{P}\{\mathbb{D}(T)\mathbb{X}(T) < 0\} = 0$, and by equivalence of $\mathbb{P}$ and $\tilde{\mathbb{P}}$ (they agree on probability-zero events), $\tilde{\mathbb{P}}\{\mathbb{D}(T)\mathbb{X}(T) < 0\} = 0$, hence
        \[\tilde{\mathbb{E}}[\mathbb{D}(T)\mathbb{X}(T)] \geq 0. \tag{B}\]
        If (A) and (B) hold simultaneously, then $\tilde{\mathbb{P}}\{\mathbb{D}(T)\mathbb{X}(T) > 0\} = 0$, and again by equivalence $\mathbb{P}\{\mathbb{X}(T) > 0\} = 0$. This contradicts requirement (c) for arbitrage. Hence no arbitrage exists.
    \end{proof}

    \item \textbf{Complete markets.} A (model of a) financial market is \textit{complete} if every derivative security can be hedged: a short position on the security can be hedged, i.e.\ every derivative security can be replicated by a portfolio made up of the $m$ stocks and the money market account.

    \item \textbf{Second Fundamental Theorem of Asset Pricing:} A market model is complete if and only if there is a \textit{unique} risk-neutral measure (i.e.\ a unique solution to the market price of risk equations).

    \item \textbf{Martingale Representation Theorem:} (see Shreve II for the actual statement). If we have a martingale w.r.t.\ the risk-neutral measure $\tilde{\mathbb{P}}$, then it can always be represented as an Ito integral. In particular, for the discounted wealth,
    \[d(\mathbb{D}(t)\mathbb{X}(t)) = \sum_{j=1}^{d}\tilde{\Pi}_j(t)\, d\tilde{\mathbb{W}}_j(t),\]
    where the $\tilde{\Pi}_j$ are adapted processes.

    \item \textbf{Why uniqueness is equivalent to completeness.} For a replicating portfolio, $d(\mathbb{D}\mathbb{X}) = \sum_I\Delta_I\, d(\mathbb{D}\mathbb{S}_I)$, so
    \[d(\mathbb{D}\mathbb{X}) = \sum_{I=1}^{m}\Delta_I\mathbb{D}\mathbb{S}_I\sum_{j=1}^{d}\sigma_{Ij}\, d\tilde{\mathbb{W}}_j = \mathbb{D}\sum_{j=1}^{d}\left(\sum_{I=1}^{m}\Delta_I\mathbb{S}_I\sigma_{Ij}\right)d\tilde{\mathbb{W}}_j.\]
    Comparing with the martingale representation,
    \[\boxed{\tilde{\Pi}_j(t) = \mathbb{D}(t)\sum_{I=1}^{m}\Delta_I(t)\mathbb{S}_I(t)\sigma_{Ij}(t), \qquad j = 1, \dots, d}\]
    These are $d$ equations in the $m$ unknowns $\Delta_I(t)$. The market is complete iff these have a solution for every choice of $\tilde{\Pi}_j$ (i.e.\ every choice of portfolio/derivative), and this happens if and only if the market price of risk equations have a unique solution.

    \item \textbf{Remark (linear algebra view).} Both conditions are statements about the rank of the $m \times d$ volatility matrix $\sigma$. $(\ast)$ has a unique solution $\vec{\Theta}$ (given that one exists) iff $\sigma$ has $d$ independent columns, and the hedging equations are solvable for every $\tilde{\Pi}$ iff $\sigma^T$ is onto $\mathbb{R}^d$, which is again $\mathrm{rank}(\sigma) = d$. Intuitively: no arbitrage needs the stocks' risk premia to be consistent (solvability of $(\ast)$), and completeness needs at least as many independent traded stocks as there are independent sources of randomness ($m \geq d$ with full rank).

\end{enumerate}
